\chapter{Results}
\label{ch:results}

This chapter reports the single-leg CFD results and what follows from them for the whole robot. It first
develops the drag-type stroke (\cref{sec:res-drag}) and characterizes the fluke (\cref{sec:res-lift}), then
compares both across speed (\cref{sec:res-compare}), improves the fluke by scheduling its pitch with speed
(\cref{sec:res-l3}), predicts the cruising speed of the robot (\cref{sec:res-selfprop}) and compares it with a first hardware
test (\cref{sec:res-hw}). All forces and powers are
per leg and for the propulsor alone unless stated otherwise.

\section{Design of the Drag-Type Stroke}
\label{sec:res-drag}

The drag-type stroke was optimized at $\U=0$ one factor at a time. The quasi-steady model of \cref{sec:fund-paddle} served as a
surrogate that proposed each change, and CFD evaluated it; the fold timing was then optimized over both switching phases on the
two-state composite. \Cref{tab:drag-design} lists the steps.

\begin{table}[t]
  \centering
  \caption[Design steps of the drag-type stroke]{Design steps of the drag-type stroke at $\U=0$. Each step changes
    one aspect of the previous one. $\Tm$ is the cycle-mean thrust of the two-state composite. All steps before the last
    switch the fan at the reversal points of the hip motion, where the added-mass correction \eqref{eq:addedmass} vanishes
    at $\U=0$; for the last step the range of the correction is given.}
  \label{tab:drag-design}
  \small
  \begin{tabularx}{\textwidth}{@{}l>{\raggedright\arraybackslash}Xcccc@{}}
    \toprule
    Stroke & Change & $T$ (s) & DUTY & $\Tm$ (mN) & $\Pm$ (mW) \\
    \midrule
    Robot gait & baseline, 30\,mm folded, \SI{60}{\degree} sinusoidal sweep & 1.25 & 0.50 & 23.4 & 6.7 \\
    V1 & shorter period, larger thrust & 0.89 & 0.49 & 51.2 & -- \\
    V2 & shorter power stroke & -- & 0.35 & 34.2 & -- \\
    V3 & sweep centred on the vertical (\SIrange{120}{70}{\degree}), slow curl & 0.81 & 0.45 & 49.3 & 16.2 \\
    V3c12 & fan folded to 12\,mm instead of 30\,mm & 0.81 & 0.45 & 56.7 & 13.7 \\
    V3c06 & folded and feathered, 5.6\,mm frontal width & 0.81 & 0.45 & 58.6 & 13.1 \\
    V3trap & V3c12 with trapezoidal velocity law (\SI{45}{\milli\second} ramps) & 0.72 & 0.382 & 75.0 & 20.8 \\
    V3trap, timed & fan closed at the start of the deceleration, $\tau_c=0.33$ & 0.72 & 0.382 & \textbf{73--83} & 21.6 \\
    \bottomrule
  \end{tabularx}
\end{table}

\paragraph{Direction of the stroke.} Centring the sweep on the vertical makes the paddle move nearly straight
backwards with its face normal to the motion. From V1 to V3 the inclination of the power-stroke force from the
swimming direction fell from \SI{19.5}{\degree} to \SI{7.9}{\degree}, and the ratio of mean vertical to mean
horizontal force from 0.52 to 0.28. V2 shows that shortening the power stroke alone loses thrust when the
recovery cannot become correspondingly cheaper.

\paragraph{Folded width.} Folding the fan to \SI{12}{\milli\metre} instead of \SI{30}{\milli\metre} reduced the
recovery drag from \SI{-10.7}{} to \SI{-3.2}{\milli\newton} and raised the mean thrust by 15\,\%. The effective
area ratio $a$, defined as the ratio of the mean force of the closed-only to the open-only run, fell from 0.48 to
0.11. With $a=0.11$, \eqref{eq:duty-opt} gives $x^*=2.2$ and $\mathrm{DUTY}^*=0.69$: the recovery should be much
faster than the power stroke, which the servo speed does not allow. The DUTY of the final stroke, 0.382, was therefore not taken from \eqref{eq:duty-opt} but found in the CFD optimization within the joint speeds that the servos allow (\cref{tab:actuators}). Feathering the folded fan to a frontal width of \SI{5.6}{\milli\metre} added only
3\,\%: the recovery drag is proportional to the width, but at \SI{12}{\milli\metre} it is only about 6\,\% of the
mean thrust. A feathering mechanism is therefore not worthwhile.

\paragraph{Velocity law.} Replacing the cosine hip velocity by a trapezoidal law with \SI{45}{\milli\second} ramps
at the same peak rate and sweep increased the power-stroke impulse by 18\,\%, as predicted by the quasi-steady
$\int v^2\,\mathrm dt$, and the mean thrust by 32\,\%, because the period also shortened from \SI{0.81}{} to
\SI{0.72}{\second}. The price is a higher peak hip moment (\SI{35.6}{} instead of \SI{19.3}{\milli\newton\metre},
$\times1.84$), a higher peak force and 52\,\% more power; the thrust per power fell from \SI{4.1}{} to
\SI{3.6}{\newton\per\watt}. The effective normal-force coefficient of the power stroke in the CFD is 2--3 times the
quasi-steady value and is highest at the start of the stroke, as expected for a starting plate \cite{Kim2011}; it
is the same for both velocity laws.

\paragraph{Fold timing.} The two-state composite (\cref{sec:cfd-composite}) allows the fold timing to be scanned
after the simulation. Without correction, the best pair is $\tau_c=0.33$, $\tau_o=-0.06$ with \SI{87.2}{\milli\newton}. The gain
from opening before the reversal, however, is the artefact described in \cref{sec:cfd-composite}, and the fan is therefore
opened at the reversal, $\tau_o=0$. Closing the fan early, at the start of the power-stroke deceleration ($\tau_c=0.33$,
0.86\,DUTY), raises the uncorrected composite from \SI{75.0}{} to \SI{82.7}{\milli\newton}: the decelerating open fan is no longer
pulled backwards by the water it has set in motion. Whether that water keeps its momentum depends on how much of it is shed as
a stopping vortex, so the added-mass correction \eqref{eq:addedmass} gives a range of \SIrange{73}{83}{\milli\newton}. At rest,
early closing is therefore worth between $-2$ and $+8$\,mN. At forward speed it is never worse and up to \SI{10}{\milli\newton} better, because after $\tau\approx0.33$ the paddle moves backwards more slowly than the oncoming flow and an open fan would only
produce drag; at \SI{0.223}{\metre\per\second}, closing at $\tau_c=0.33$ gives \SIrange{24}{26}{\milli\newton}, against \SIrange{14}{26}{\milli\newton} for closing at the reversal. The final stroke, called V3trap in the rest of this thesis (\cref{fig:composite,fig:kinematics}a),
produces \SIrange{73}{83}{\milli\newton} per leg at rest, 3.1--3.5 times the robot's current gait, at \SI{21.6}{\milli\watt} and
\SIrange{3.4}{3.8}{\newton\per\watt}. These rules agree with those found for crustacean and webbed-foot propulsors: open fully at
the start of the power stroke, close early, keep transitions short \cite{Santos2026,Lin2024}.

\section{The Lift-Type Fluke with a Fixed Pitch Schedule}
\label{sec:res-lift}

\Cref{tab:lift} lists the fluke results with the L0 motion at five speeds. At rest the inflow is purely vertical,
the fluke is stalled ($\alpha_{75}=\SI{73}{\degree}$) and acts as a small paddle; it produces
\SI{45.5}{\milli\newton} at \SI{18.1}{\milli\watt}, about twice the thrust of the robot gait with the fan paddle.
With increasing speed the inflow angle and the angle of attack fall, and the thrust decreases almost linearly, to
\SI{12.7}{\milli\newton} at \SI{0.30}{\metre\per\second}. The input power stays between \SI{15.6}{} and
\SI{18.1}{\milli\watt}. The efficiency rises to a maximum of 0.29 at \SI{0.223}{\metre\per\second}, the speed for
which L0 was designed ($\mathit{St}\approx0.35$ based on the pin excursion, $\alpha_{75}=\SI{19.9}{\degree}$). At
\SI{0.30}{\metre\per\second} the mid-stroke angle of attack drops to about \SI{10}{\degree}, the efficiency falls
to 0.245, and a mean vertical force of \SI{-29}{\milli\newton} appears, more than twice the thrust.

\begin{table}[t]
  \centering
  \caption[Fluke with the L0 motion]{Fluke with the L0 motion (fixed pitch schedule) at different speeds.
    $\mathit{St}$ is based on the peak-to-peak pin excursion of \SI{39.2}{\milli\metre}.}
  \label{tab:lift}
  \small
  \begin{tabular}{@{}lccccc@{}}
    \toprule
    $\U$ (\si{\metre\per\second}) & 0 & 0.08 & 0.15 & 0.223 & 0.30 \\
    \midrule
    $\Tm$ (mN) & 45.5 & 36.5 & 29.1 & 21.4 & 12.7 \\
    $\Pm$ (mW) & 18.1 & 17.1 & 16.6 & 16.3 & 15.6 \\
    $\etaF$ & -- & 0.171 & 0.263 & 0.293 & 0.245 \\
    $\Tm/\Pm$ (\si{\newton\per\watt}) & 2.51 & 2.13 & 1.75 & 1.31 & 0.81 \\
    $\mathit{St}$ & -- & 0.98 & 0.52 & 0.35 & 0.26 \\
    \bottomrule
  \end{tabular}
\end{table}

The wake at \SI{0.223}{\metre\per\second} (\cref{fig:wake}) shows the expected pattern of a thrust-producing
foil: each half stroke sheds a vortex of alternating sign from the trailing edge, and the vortices form a reverse
von K\'arm\'an street that convects downstream.

\begin{figure}[t]
  \centering
  \includegraphics[width=\textwidth]{fig_wake}
  \caption[Wake of the fluke]{Spanwise vorticity in the mid-span plane of the fluke (L0,
    $\U=\SI{0.223}{\metre\per\second}$) at four phases of the third cycle, hull frame; the flow is from right to left
    relative to the fluke. Black line: fluke chord.}
  \label{fig:wake}
\end{figure}

Additional pitch schedules for target angles of attack of \SI{10}{\degree} and \SI{28}{\degree} at \SI{0.223}{\metre\per\second}
produce more thrust than L0 but at lower efficiency (0.20 and 0.24, \cref{tab:app-fluke}); among the cases at this speed the
efficiency is highest for $\alpha_{75}\approx\SI{20}{\degree}$, in line with the efficient range of flapping foils \cite{Read2003}.

\section{Drag versus Lift across Speed}
\label{sec:res-compare}

\Cref{fig:thrust-speed} compares thrust, power and efficiency of the propulsors from rest to
\SI{0.40}{\metre\per\second}; \cref{tab:compare} gives the values.

\begin{figure}[t]
  \centering
  \includegraphics[width=\textwidth]{fig_thrust_speed}
  \caption[Thrust, power and efficiency versus speed]{Single-leg (a) thrust, (b) input power and (c) Froude
    efficiency of the propulsors versus speed. Shaded: range of the added-mass correction of the paddle composites
    \eqref{eq:addedmass}; dotted: V3trap composite without correction.}
  \label{fig:thrust-speed}
\end{figure}

\begin{table}[t]
  \centering
  \caption[Propulsor comparison across speed]{Single-leg thrust $\Tm$ (mN) / input power $\Pm$ (mW) / Froude
    efficiency $\etaF$ across speed. Paddle values include the added-mass correction \eqref{eq:addedmass} as a range; the
    uncorrected composite of V3trap gives 82.7, 75.6, 58.2, 35.7 and \SI{11.7}{\milli\newton}.}
  \label{tab:compare}
  \footnotesize
  \begin{tabular}{@{}ccccc@{}}
    \toprule
    $\U$ (\si{\metre\per\second}) & Paddle, robot gait & Paddle, V3trap & Fluke, fixed (L0/L2) & Fluke, scheduled (L3) \\
    \midrule
    0     & 23.4 / 6.7 / --         & \textbf{73--83} / 21.6 / --                & 45.5 / 18.1 / --      & -- \\
    0.08  & 4--5 / 5.8 / 0.05--0.07 & \textbf{66--71} / 21.9 / 0.24--0.26        & 36.5 / 17.1 / 0.17    & -- \\
    0.15  & $<0$ / 6.5 / --         & \textbf{48--50} / 20.7 / \textbf{0.35--0.37} & 29.1 / 16.6 / 0.26  & -- \\
    0.223 & --                      & 24--26 / 19.2 / 0.28--0.30               & 21.4 / 16.3 / 0.29    & -- \\
    0.30  & --                      & $-4$--2 / 17.8 / $\approx0$                & 12.7 / 15.6 / 0.24    & \textbf{13.9} / 10.5 / \textbf{0.40} \\
    0.40  & --                      & --                                       & --                    & 10.5 / 11.4 / 0.37 \\
    \bottomrule
  \end{tabular}
\end{table}

\paragraph{Low speed.} Up to about \SI{0.15}{\metre\per\second} the optimized paddle is clearly superior. It produces
1.6--2.0 times the thrust of the fluke at 19--28\,\% more power, and at \SI{0.15}{\metre\per\second} it reaches an efficiency of
0.35--0.37, higher than the fluke's maximum of 0.29. Per area there is no difference at rest: the paddle produces
\SIrange{3.9}{4.5}{} and the fluke \SI{4.1}{\milli\newton\per\square\centi\metre}. At \SI{0.223}{\metre\per\second} the fluke
already produces more per area (\SI{1.9}{} versus \SIrange{1.3}{1.4}{\milli\newton\per\square\centi\metre}). The robot gait, in
contrast, loses its thrust almost immediately: at \SI{0.08}{\metre\per\second} only \SIrange{4}{5}{\milli\newton} remain, and at
\SI{0.15}{\metre\per\second} the net force is a drag.

\paragraph{Speed limit of the paddle.} The paddle's thrust falls steeply with speed and vanishes at about
\SI{0.30}{\metre\per\second}. The split into power stroke and recovery (\cref{fig:drag-split}, values in \cref{tab:app-split}) shows why. The power-stroke
thrust of V3trap decreases by 42\,\%, from \SI{88.0}{} to \SI{50.9}{\milli\newton}, because the relative velocity during the power
stroke falls. The recovery drag, however, grows more than sevenfold, from \SI{-5.3}{} to \SI{-39.1}{\milli\newton}, because the
folded paddle and the rod now move forward against the oncoming flow. At speed, opening the fan also costs the momentum
$\Delta m_a\U$ per cycle, \SIrange{12}{16}{\milli\newton} at \SIrange{0.223}{0.30}{\metre\per\second}. At
\SI{0.30}{\metre\per\second} the recovery and the opening together consume the whole power-stroke thrust. For the robot gait, with its
wider folded fan and slower stroke, the recovery drag exceeds the power-stroke thrust already at about
\SI{0.1}{\metre\per\second}. The speed limit of a drag-type paddle on this leg is therefore set by what the paddle has to do outside
the power stroke, not by the speed of the paddle in the power stroke.

\begin{figure}[t]
  \centering
  \includegraphics[width=\textwidth]{fig_drag_split}
  \caption[Power stroke and recovery of the paddle]{Contributions of power stroke and recovery to the mean thrust of
    the fan paddle. (a) Absolute values for V3trap and the robot gait; shaded: net thrust including the added-mass correction
    \eqref{eq:addedmass}. (b) Magnitudes relative to the power-stroke thrust at $\U=0$.}
  \label{fig:drag-split}
\end{figure}

\paragraph{Crossover.} The fluke loses thrust more slowly, because it neither recovers against the flow nor changes its shape:
its thrust at \SI{0.30}{\metre\per\second} is 28\,\% of its value at rest, whereas the paddle's is close to zero. Interpolating the
curves, the thrust of the fluke exceeds that of V3trap above \SIrange{0.23}{0.24}{\metre\per\second}, and its efficiency above
\SIrange{0.22}{0.23}{\metre\per\second}. This holds for both pitch schedules, because the speed-scheduled pitch was simulated only at \SIlist{0.30;0.40}{\metre\per\second} and shares the L0 data below. Including the discretization uncertainty of
\cref{sec:ver-drag}, the crossover lies at $0.24\pm\SI{0.02}{\metre\per\second}$ for thrust and at $0.22\pm\SI{0.02}{\metre\per\second}$
for efficiency. Without the added-mass correction, the crossovers would lie at \SI{0.295}{} and \SI{0.285}{\metre\per\second}; the
correction moves them by about \SI{0.06}{\metre\per\second} but does not change their existence.

\section{Speed-Scheduled Fluke Pitch}
\label{sec:res-l3}

With the fixed L0 motion, the fluke's angle of attack falls with speed, which explains its efficiency drop above
\SI{0.223}{\metre\per\second}. The speed-scheduled pitch \eqref{eq:l3} restores the angle of attack
(\cref{fig:l3}). At \SI{0.30}{\metre\per\second} it reduces the pitch amplitude from about \SI{40}{\degree} to
\SI{21.6}{\degree}, and the angle of attack rises accordingly: $\alpha_{75}$ increases from \SI{13.7}{\degree} to \SI{21.0}{\degree}
(\cref{tab:l3,tab:app-fluke}).

\begin{table}[t]
  \centering
  \caption[Speed-scheduled pitch]{Speed-scheduled pitch (L3) compared with the fixed L0 motion (L2). $\bar F_z$: mean vertical
    force; $|M_h|$: peak hinge moment about the pin.}
  \label{tab:l3}
  \small
  \begin{tabular}{@{}lcccccc@{}}
    \toprule
    Case & $\U$ (\si{\metre\per\second}) & $\Tm$ (mN) & $\Pm$ (mW) & $\etaF$ & $\bar F_z$ (mN) & $|M_h|$ (\si{\milli\newton\metre}) \\
    \midrule
    L2, fixed & 0.30 & 12.7 & 15.6 & 0.245 & $-29.4$ & 12.6 \\
    L3, scheduled & 0.30 & \textbf{13.9} & \textbf{10.5} & \textbf{0.398} & $-0.9$ & 7.4 \\
    L3, scheduled & 0.40 & 10.5 & \textbf{11.4} & \textbf{0.368} & $-0.9$ & 7.3 \\
    \bottomrule
  \end{tabular}
\end{table}

At \SI{0.30}{\metre\per\second} the scheduled pitch raises the thrust by 9\,\%, lowers the power by 33\,\% and raises the
efficiency from 0.245 to 0.40, while the paddle produces almost no thrust at this speed. It also removes the mean vertical force
($-29$ to $\SI{-1}{\milli\newton}$), reduces the peak hinge moment by 40\,\% and flattens the force history: the
peak-to-peak streamwise force falls from about \SI{600}{} to \SI{190}{\milli\newton}. At
\SI{0.40}{\metre\per\second} the efficiency is 0.37. Over the three designed speeds the fluke efficiency is 0.29 (L0, 0.223), 0.40 (L3,
0.30) and 0.37 (L3, 0.40), with a maximum near \SI{0.30}{\metre\per\second}. At this speed the Strouhal number based on the pin
excursion is 0.26, inside the range of 0.2--0.4 of efficient cruising in animals and close to the mean value measured for odontocetes \cite{Taylor2003,Rohr2004}; based on the trailing-edge excursion it is somewhat higher. The schedule needs no flow
information, only $\dot h_\text{max}$ and $\U$. On the hardware it corresponds to a pitch amplitude that decreases with speed; how
closely a passive spring can provide it is discussed in \cref{sec:disc-limits-prop}.

\begin{figure}[t]
  \centering
  \includegraphics[width=\textwidth]{fig_l3}
  \caption[Speed-scheduled pitch]{Fixed L0 motion versus speed-scheduled pitch. (a) Prescribed fluke pitch (hinge heave
    dotted, right axis). (b) Resulting angle of attack. (c) Streamwise force, smoothed.}
  \label{fig:l3}
\end{figure}

\section{Self-Propulsion Estimate}
\label{sec:res-selfprop}

\paragraph{Resistance.} Three resistance curves bound the resistance of the robot without its propulsors
(\cref{fig:selfprop}a):
\begin{itemize}
  \item $D_\text{low}=\SI{10.8}{\milli\newton}\,(\U/\SI{0.2}{\metre\per\second})^{1.72}$: the bare hull, from separate
        two-phase CFD of the hull at \SIlist{0.1;0.2}{\metre\per\second};
  \item $D_\text{mid}=D_\text{low}+\SI{15.2}{\milli\newton}\,(\U/\SI{0.2}{\metre\per\second})^2$: hull plus the submerged
        linkages, about \SI{26}{\milli\newton} at \SI{0.2}{\metre\per\second}; the linkage term, equivalent to
        $C_DA\approx\SI{7.6}{\square\centi\metre}$, is an engineering estimate, not a measurement; this is the best estimate;
  \item $D_\text{up}=\SI{77.6}{\milli\newton}\,(\U/\SI{0.2}{\metre\per\second})^2$: the whole robot at rest with its fan
        paddles open, from CFD; it counts the paddles' own drag, which is already contained in $\Tm$, and is therefore an
        upper bound.
\end{itemize}
The thrust curves are interpolated with monotone cubic splines and extrapolated linearly beyond the last CFD point.

\begin{figure}[t]
  \centering
  \includegraphics[width=\textwidth]{fig_selfprop}
  \caption[Self-propulsion estimate]{Self-propulsion estimate. (a) Four-leg thrust versus the resistance range; circles mark
    the cruising speeds for $D_\text{mid}$; shaded: range of the added-mass correction of the paddle. Dotted: linear extrapolation. Below \SI{0.30}{\metre\per\second}, where the speed-scheduled pitch was not simulated, its curve uses the fixed-pitch data. (b) Quasi-steady acceleration from rest for $D_\text{mid}$; bars on the right give
    the range of cruising speeds from $D_\text{up}$ to $D_\text{low}$.}
  \label{fig:selfprop}
\end{figure}

\begin{table}[t]
  \centering
  \caption[Predicted cruising speed]{Predicted cruising speed and four-leg input power for $D_\text{mid}$, with the range
    from $D_\text{up}$ to $D_\text{low}$. Paddle values are given as the range of the added-mass correction. $^*$ extrapolated beyond
    the CFD data. $t_{90}$: time to reach 90\,\% of the cruising speed from rest.}
  \label{tab:selfprop}
  \small
  \begin{tabular}{@{}lccccc@{}}
    \toprule
    Propulsor & $\U_c$ (\si{\metre\per\second}) & Range (\si{\metre\per\second}) & $4\Pm$ (mW) & $4\Pm/\U_c$ (\si{\joule\per\metre}) & $t_{90}$ (s) \\
    \midrule
    Drag, robot gait & 0.089--0.095 & 0.082--0.098 & 23.3 & 0.25--0.26 & 0.94--0.99 \\
    Drag, V3trap & 0.260--0.269 & 0.223--0.289 & 73--74 & 0.27--0.28 & 0.67--0.69 \\
    Lift, fixed pitch & 0.294 & 0.214--0.350$^*$ & 62.7 & 0.21 & 1.46 \\
    Lift, scheduled pitch & \textbf{0.299} & 0.214--0.422$^*$ & \textbf{42.0} & \textbf{0.14} & 1.57 \\
    \bottomrule
  \end{tabular}
\end{table}

\paragraph{Cruising speed.} For the best-estimate resistance, the fluke propels the robot at \SI{0.30}{\metre\per\second}, or
1.3 hull lengths per second, and the optimized paddle at \SIrange{0.26}{0.27}{\metre\per\second} (\cref{tab:selfprop}). Both
operating points lie just above the crossover, where the fluke is stronger, but close to it, so the difference in speed is small.
The power differs clearly: the speed-scheduled fluke needs \SI{42}{\milli\watt}, 33\,\% less than the fixed-pitch fluke
(\SI{63}{\milli\watt}) and 43\,\% less than the paddle (\SI{73}{\milli\watt}), which is also slower. The robot's current gait reaches
only about \SI{0.09}{\metre\per\second}. The ranges are wide because the resistance is uncertain, and which propulsor is faster
depends on where the resistance curve crosses the thrust curves. With the high resistance $D_\text{up}$, the operating point
moves below the crossover and the paddle is slightly faster (\SIrange{0.22}{0.23}{} versus \SI{0.21}{\metre\per\second}); with the
low resistance $D_\text{low}$, it moves further above the crossover and the fluke is clearly faster
(\SIrange{0.35}{0.42}{\metre\per\second}, extrapolated, versus \SIrange{0.28}{0.29}{\metre\per\second}).

\paragraph{Acceleration.} The acceleration from rest is integrated with the measured mass of \SI{0.419}{\kilo\gram} plus 10\,\% added
mass; the mass affects only the time scale, not the cruising speed. The paddle accelerates the robot faster: it reaches 90\,\% of its cruising speed after
\SI{0.68}{\second}, the fluke after \SIrange{1.46}{1.57}{\second} (\cref{fig:selfprop}b). This reflects the paddle's higher thrust at low
speed and matches the observation that rowing appendages favour acceleration and flapping ones cruising
\cite{Walker2000}.

\section{First Hardware Test}
\label{sec:res-hw}

The robot with flukes (\cref{sec:plat-hw}) was timed in a \SI{1}{\metre} pool: it started from rest with its stern against one
wall and the time was stopped when its bow hit the opposite wall, a distance of \SI{0.738}{\metre} (pool length minus the overall
length of the hull, \SI{262}{\milli\metre}). The stroke frequency was set to \SI{2}{} and \SI{4}{\hertz} in the firmware; there was
one hand-timed run for each setting. The same run was simulated with the quasi-steady model of \cref{fig:selfprop}b, the L0 thrust at
\SI{2}{\hertz} and the measured mass (\cref{tab:hw}).

\begin{table}[htb]
  \centering
  \caption[First hardware test]{First hardware test: time and mean speed over \SI{0.738}{\metre} from rest. Frequencies are the
    values set in the firmware. Model: L0 thrust at \SI{2}{\hertz} with the three resistance curves of \cref{sec:res-selfprop}.}
  \label{tab:hw}
  \small
  \begin{tabular}{@{}lccc@{}}
    \toprule
    & Frequency set (Hz) & Time (s) & Mean speed (\si{\metre\per\second}) \\
    \midrule
    Hardware & 2 & 5.2 & 0.142 \\
    Hardware & 4 & 2.8 & 0.264 \\
    Model, $D_\text{low}$ / $D_\text{mid}$ / $D_\text{up}$ & 2 & 2.9 / 3.2 / 3.9 & 0.25 / 0.23 / 0.19 \\
    \bottomrule
  \end{tabular}
\end{table}

\paragraph{Thrust.} At \SI{2}{\hertz}, where the stroke is close to L0, the hardware needs 1.3 to 1.6 times as long as predicted with $D_\text{mid}$ to $D_\text{up}$ (1.8 times with $D_\text{low}$, which does not apply to the hardware).
The resistance alone cannot explain this: with the CFD thrust, a resistance of \SI{205}{\milli\newton} at
\SI{0.2}{\metre\per\second} would be required, 2.6 times $D_\text{up}$. This corresponds to $C_DA\approx\SI{103}{\square\centi\metre}$,
or $C_D\approx4$ on the submerged frontal area of the hull of \SI{24}{\square\centi\metre}, far above any bluff body. $D_\text{up}$
itself corresponds to $C_D\approx1.6$ and is a generous bound even for the hull moving blunt end first, while $D_\text{low}$, the bare
hull moving pointed end first, does not apply to the hardware. Between $D_\text{mid}$ and $D_\text{up}$, the run is reproduced with
27--47\,\% of the CFD thrust, roughly a quarter to a half, and a steady speed of about \SIrange{0.16}{0.20}{\metre\per\second} at
\SI{2}{\hertz}. Part of the deficit is expected from the stroke itself: the piecewise-linear trajectory has a peak heave speed about
15\,\% below that of L0 (\cref{sec:plat-hw}), which costs of the order of 10--20\,\% of the thrust. The remaining factor of about two
has several possible causes that the test cannot separate: the passive pitch, whose spring stiffness has not been verified against the
angle of attack of L0; ventilation where the fluke breaks the surface; the reversed hull; and interaction between the legs. The
self-propulsion speeds of \cref{tab:selfprop} are therefore upper bounds for the present hardware.

\paragraph{Frequency.} The ratio of the two runs does not depend on the thrust level. If the resistance grows with $\U^2$ and the
thrust scales with the square of the stroke velocity at a fixed Strouhal number, the time needed to cover a given distance from
rest is inversely proportional to the product of frequency and amplitude. The measured ratio of \SI{5.2}{} to \SI{2.8}{\second},
1.86, is close to the value of 2 expected if the servos reach the full amplitude at \SI{4}{\hertz}, although the peak servo speed of about \SI{800}{\degree\per\second}, estimated from the firmware trajectory, is close to their no-load speed. A more precise statement is not possible: a timing error of
\SI{0.2}{\second} per run moves the ratio between 1.67 and 2.08, and the passive pitch changes with frequency, because the
hydrodynamic moment grows with the square of the stroke velocity while the spring moment does not. Taken at face value, the same
scaling puts the steady speed at \SI{4}{\hertz} at roughly \SIrange{0.3}{0.4}{\metre\per\second}, above the CFD-based estimate at
\SI{2}{\hertz}, for about six times the hydrodynamic power.
